\chapter{Второй выход}
% полная версия: chapters/16_chemoutput.tex

У машины два выхода. Быстрый --- мышцы. Медленный --- химия тела: мозг управляет сердцем, сосудами, железами и через кровь рассылает сообщения всему телу сразу. Кровь --- самая широковещательная шина из всех: одно сообщение доходит до каждой клетки примерно за минуту, а смысл ему, как всегда, придаёт «замок» получателя.

\section*{Газ и тормоз}

Сердце бьётся само, у него свой метроном. Мозг лишь подкручивает темп двумя проводами противоположного знака. \nt{НОРА} --- педаль газа: разгоняет, но медленно. \nt{АЦЕТ} --- тормоз: замедляет быстро, потому что ему не нужна длинная цепочка химических реакций. И тут наконец находится работа для типа \nt{АДРЕ}: адреналин, выброшенный в кровь, --- та же педаль газа, только нажатая сразу для всего тела.

\pencilfig{16}{\draw[penlite=pcGreen, wash=pcGreen, rounded corners=2pt] (0,0) rectangle (0.45,1.1);
\draw[penlite=pcRed, wash=pcRed, rounded corners=2pt] (0.85,0) rectangle (1.65,0.65);
\node[lbl, anchor=north] at (0.22,-0.05) {газ}; \node[lbl, anchor=north] at (1.25,-0.05) {тормоз};
\draw[pen=pcRed, wash=pcRed] (3.2,0.15) .. controls (2.5,0.65) and (2.75,1.3) .. (3.2,1.0) .. controls (3.65,1.3) and (3.9,0.65) .. (3.2,0.15) -- cycle;
\draw[pcGray, dashed, line width=0.5pt, ->] (0.5,0.9) -- (2.6,0.9); \draw[pcGray, dashed, line width=0.5pt, ->] (1.7,0.45) -- (2.65,0.55);}{Сердце ведут два провода противоположного знака: газ и тормоз.}

\section*{Каскад с обратной связью}

Многие гормоны управляются каскадом: мозг приказывает первой железе, та --- второй, вторая выделяет гормон, а гормон сообщает мозгу, что его уже достаточно. Это термостат из трёх ступеней. Инженер знает, чем опасна такая схема: если петлю сделать слишком «чувствительной», она начнёт раскачиваться --- уровень будет не держаться, а колебаться. Поэтому точность такого регулятора ограничена: сделать его одновременно устойчивым и очень точным нельзя.

\section*{Импульсы, а не уровень}

Некоторые гормоны работают только импульсами. Если подавать их постоянно, железа-получатель «устаёт» и отключается; если короткими импульсами раз в час --- работает. Получатель читает не уровень, а ритм, --- как синапс из главы про азбуку Морзе, который слышит только перемены.

\section*{Сутки}

И наконец, в теле есть медленный генератор с периодом около суток. Его естественный период чуть отличается от двадцати четырёх часов, и каждое утро свет подстраивает его, как радиоприёмник подстраивается под станцию. Не путайте его с тактовым генератором компьютера: он не отмеряет шаги вычисления. Он переключает режимы работы всей машины --- день и ночь, бодрствование и сон.

\fullversion{16}
